% !TeX root = ../../Book2.tex
\section{卷积}
\label{b2:sec:0303}

频率上的乘法对应空间中的卷积。要使用这个对应关系，先须证明卷积本身存在，并明确它落在哪个函数空间。可积性在这里不仅是形式上的标签，也是交换积分的依据。

\subsection{卷积代数}
\label{b2:subsec:030301}

\begin{definition}\label{b2:def:03-convolution}
对可测函数 \(f,g\)，若积分绝对收敛，定义它们在 \(x\) 处的\term[juanji]{卷积}[convolution]为
\[
 (f*g)(x)\coloneqq\int_{\mathbb R^d}f(x-y)\cdot g(y)\,\mathrm{d}y.
\]
\end{definition}
\symidx{\(f*g\)：欧氏空间卷积}

\begin{proposition}\label{b2:prop:03-l1-algebra}
若 \(f,g\in L^1(\mathbb R^d)\)，则卷积几乎处处存在，属于 \(L^1\)，且
\[
 \|f*g\|_1\leq\|f\|_1\cdot\|g\|_1.
\]
卷积在 \(L^1\) 等价类上良定义，满足交换律、结合律和双线性。
\end{proposition}
\begin{proof}
可测函数可选取 Borel 代表；函数 \((x,y)\mapsto f(x-y)g(y)\) 因而联合可测。对绝对值用 Tonelli 和平移不变性，
\[
 \int\!\int|f(x-y)|\cdot|g(y)|\,\mathrm{d}y\,\mathrm{d}x
 =\|f\|_1\cdot\|g\|_1<\infty.
\]
这既保证几乎处处绝对收敛，又给出范数界。改变任一代表所产生的差，按同一估计为零，故等价类良定义。令 \(z=x-y\) 得交换律。三个函数的绝对值重积分等于三个范数之积，所以可用 Fubini 重排；两种卷积括号都等于
\[
 \int\!\int f(x-y-z)\cdot g(y)h(z)\,\mathrm{d}y\,\mathrm{d}z
\]
几乎处处。线性来自积分的线性。
\end{proof}

结合 \(L^1\) 的完备性，这使 \(L^1\) 成为卷积运算下的 Banach 代数。它没有单位元：若 \(e*f=f\) 对所有 \(f\in L^1\) 成立，后面将由变换证明 \(\widehat e\) 必须恒等于 \(1\)，与 Riemann--Lebesgue 引理矛盾；在分布空间中才会出现相应的 Dirac 单位。

\subsection{Young 不等式}
\label{b2:subsec:030302}

\begin{theorem}[Young，卷积形式]\label{b2:thm:03-young}
设 \(1\leq p,q,r\leq\infty\) 且 \(1/p+1/q=1+1/r\)，其中 \(1/\infty=0\)。若 \(f\in L^p\)、\(g\in L^q\)，则卷积几乎处处存在，并且
\[
 \|f*g\|_r\leq\|f\|_p\cdot\|g\|_q.
\]
\end{theorem}
\begin{proof}
若 \(r=\infty\)，指数 \(p,q\) 共轭，对每个 \(x\) 使用 Hölder 即得积分的绝对收敛与一致上界。

若 \(r<\infty\)，则 \(p,q\leq r\)。先对有界紧支集函数计算，把绝对值乘积拆成
\[
 |f(x-y)|\cdot|g(y)|
 =\bigl(|f(x-y)|^p\cdot|g(y)|^q\bigr)^{1/r}
    |f(x-y)|^{1-p/r}\cdot|g(y)|^{1-q/r}.
\]
Hölder 的三个倒指数为 \(1/r\)、\(1/p-1/r\)、\(1/q-1/r\)，和为 \(1\)。某个倒指数为零时相应因子为 \(1\)，将其省去。于是
\[
 |f*g(x)|^r
 \leq\left(\int |f(x-y)|^p\cdot|g(y)|^q\,\mathrm{d}y\right)
       \|f\|_p^{r-p}\cdot\|g\|_q^{r-q}.
\]
对 \(x\) 积分，Tonelli 给出右侧为 \(\|f\|_p^r\cdot\|g\|_q^r\)。一般函数先对绝对值作幅值和区域的递增截断，所得非负卷积由单调收敛及该统一范数界确定一个有限的 \(L^r\) 函数。这保证原复值积分几乎处处绝对收敛，积分三角不等式再给出所需估计。零范数的情形直接成立。
\end{proof}

这称为\term[young-juanji-budengshi]{Young 卷积不等式}[Young's convolution inequality]，与标量乘积的 Young 不等式区别在于对象和指数关系。常用情形是 \(L^1*L^p\subset L^p\)，以及 \(L^p*L^{p'}\subset L^\infty\)。后一种情形还可利用平移连续性推出连续代表。

\subsection{平移与卷积}
\label{b2:subsec:030303}

第一卷引理 13.4.5 已证明：若 \(1\leq p<\infty\)，则 \(\|\tau_hf-f\|_p\to0\)。该结论首先对连续紧支集函数成立，再由平移等距性和稠密性推广。无穷端点不在其中；区间示性函数的小平移在 \(L^\infty\) 中始终保持误差 \(1\)。

\begin{proposition}\label{b2:prop:03-convolution-continuity}
若 \(1<p<\infty\)、\(f\in L^p\)、\(g\in L^{p'}\)，则卷积具有有界一致连续的代表。\(L^1*L^\infty\) 的卷积也有这一性质。
\end{proposition}
\begin{proof}
Hölder 保证每个点的积分绝对收敛。由变量替换，
\[
 \|\tau_h(f*g)-f*g\|_\infty
 \leq\|\tau_hf-f\|_p\cdot\|g\|_{p'}\longrightarrow0.
\]
在 \(L^1*L^\infty\) 情形，将平移放在 \(L^1\) 因子上即可。
\end{proof}

连续性在这里由两个因子的整体平移行为取得，并非因为原函数本身连续。另一方面，\(L^1*L^1\) 一般只有可积性，不能仅凭卷积符号就称所得函数光滑。若某一因子为光滑紧支集函数，才可将导数移到该因子，并由可积控制证明结果光滑。


\begin{corollary}\label{b2:cor:03-convolution-c0}
若 \(1<p<\infty\)、\(f\in L^p\)、\(g\in L^{p'}\)，则 \(f*g\in C_0(\mathbb R^d)\)。
\end{corollary}
\begin{proof}
取连续紧支集函数 \(f_n\to f\) 于 \(L^p\)、\(g_n\to g\) 于 \(L^{p'}\)。可取两列范数有界。由 Hölder 型卷积界，
\[
 \|f*g-f_n*g_n\|_\infty
 \leq\|f-f_n\|_p\cdot\|g\|_{p'}
 +\|f_n\|_p\cdot\|g-g_n\|_{p'}\longrightarrow0.
\]
每个 \(f_n*g_n\) 连续且支撑在两个紧支集的和中，所以属于 \(C_0\)。一致极限仍属于 \(C_0\)。
\end{proof}

两个指数都有限，使两边都能用紧支集函数逼近。端点 \(L^1*L^\infty\) 不保证在无穷远消失：取 \(g=1\)，则卷积恒为 \(\int f\)。这与已经证明的有界一致连续性不矛盾。卷积的“平滑作用”应拆成可积性、连续性、无穷远行为和可微性分别判断，不能把这些性质合成一个没有条件的结论。

\subsection{Fourier 变换与卷积}
\label{b2:subsec:030304}

\begin{theorem}[卷积定理]\label{b2:thm:03-convolution-transform}
若 \(f,g\in L^1(\mathbb R^d)\)，则
\[
 \widehat{f*g}(\xi)=\widehat f(\xi)\cdot\widehat g(\xi).
\]
\end{theorem}
\begin{proof}
定义积分中的绝对值重积分等于 \(\|f\|_1\cdot\|g\|_1\)，故可交换次序。令 \(z=x-y\)，指数分解为 \(e^{-2\pi\mathrm{i}z\cdot\xi}e^{-2\pi\mathrm{i}y\cdot\xi}\)，两个积分分别为 \(\widehat f\) 与 \(\widehat g\)。
\end{proof}

这个\term[juanji-dingli]{卷积定理}[convolution theorem]将叠加平移的操作对角化。现在也可补足单位元不存在的论证：若 \(e*f=f\)，固定频率 \(\xi\) 后取单位质量非负函数的调制，使 \(\widehat f(\xi)=1\)，便有 \(\widehat e(\xi)=1\)。这与 \(\widehat e\in C_0\) 不相容。下一节改用一族核逐渐接近单位作用。
